\chapter{Detect, Calibrate, Constrain, Shield, and Certify}
\label{ch:dc3s-kernel}

The five-stage ORIUS kernel gives the monograph its system spine.

\section{Detect}
The detection stage computes observation reliability from freshness, missingness,
timing, spikes, and consistency features.  It is the first point at which the runtime
admits that the observation channel may no longer deserve full trust.

\section{Calibrate}
The calibration stage translates that loss of trust into a wider uncertainty surface.
In the current ORIUS implementation this is achieved through conformal-style or
reliability-aware predictive inflation, but the chapter treats the operation more
generally: the state set consistent with observation must expand when observation
quality degrades.

\section{Constrain}
The constrain stage converts the state set into an admissible action set.  This is
where ORIUS differs from alarm-only monitoring.  It is not enough to know that
telemetry is degraded; the admissible action set must be recomputed accordingly.

\section{Shield}
The shield stage repairs the candidate action into one that is admissible under the
tightened set or replaces it with a fallback action when no safe action remains.

\section{Certify}
The certificate stage records why the action was permitted, what the observation
quality was, what uncertainty and margin were assumed, and what lifecycle state the
runtime is now in.  This stage is what lets ORIUS move from control logic to a real
governance story.
